\documentclass[UTF8,openany,zihao=-4,fontset=none]{ctexbook}
\newcommand{\ManualTitle}{企业端软件使用手册}
\newcommand{\ManualShortTitle}{企业端软件使用手册}
\usepackage[a4paper,top=25mm,bottom=24mm,left=25mm,right=22mm,headheight=15pt]{geometry}
\usepackage{fontspec}
\usepackage{xeCJK}
\setmainfont{DejaVu Serif}
\setsansfont{DejaVu Sans}
\setmonofont{DejaVu Sans Mono}
\setCJKmainfont[AutoFakeBold=2.0]{Droid Sans Fallback}
\setCJKsansfont[AutoFakeBold=2.0]{Droid Sans Fallback}
\setCJKmonofont{Droid Sans Fallback}

\usepackage{booktabs}
\usepackage{amssymb}
\usepackage{longtable}
\usepackage{tabularx}
\usepackage{array}
\usepackage{enumitem}
\usepackage{xcolor}
\usepackage{fancyhdr}
\usepackage{hyperref}
\usepackage[most]{tcolorbox}
\usepackage{titlesec}
\usepackage{lastpage}
\usepackage{microtype}

\definecolor{ManualBlue}{HTML}{17365D}
\definecolor{ManualTeal}{HTML}{167D8D}
\definecolor{ManualLight}{HTML}{EDF4F7}
\definecolor{ManualOrange}{HTML}{B85C00}
\definecolor{ManualRed}{HTML}{A61B1B}
\definecolor{ManualGray}{HTML}{53606D}
\definecolor{ManualLine}{HTML}{C9D5DD}

\hypersetup{
  colorlinks=true,
  linkcolor=ManualBlue,
  urlcolor=ManualTeal,
  pdftitle={\ManualTitle},
  pdfauthor={MineGuard 项目组},
  pdfsubject={软件使用手册},
  pdfkeywords={MineGuard, 煤炭, 监管, 企业填报, 使用手册},
  bookmarksnumbered=true,
  bookmarksopen=true
}

\ctexset{
  chapter={
    format=\Huge\bfseries\color{ManualBlue},
    name={第,章},
    number=\chinese{chapter},
    beforeskip=0pt,
    afterskip=22pt
  },
  section={format=\Large\bfseries\color{ManualBlue}},
  subsection={format=\large\bfseries\color{ManualTeal}}
}
\titleformat{\paragraph}[runin]{\bfseries\color{ManualBlue}}{}{0pt}{}[。]

\setlength{\parindent}{2em}
\setlength{\parskip}{0.35em}
\setlist{itemsep=0.25em,topsep=0.35em,leftmargin=2.2em}
\setlist[enumerate,1]{label=\arabic*.}
\renewcommand{\arraystretch}{1.35}
\newcolumntype{Y}{>{\raggedright\arraybackslash}X}
\newcolumntype{P}[1]{>{\raggedright\arraybackslash}p{#1}}

\pagestyle{fancy}
\fancyhf{}
\fancyhead[L]{\small\color{ManualGray}\ManualShortTitle}
\fancyhead[R]{\small\color{ManualGray}文档版本 V1.0}
\fancyfoot[L]{\small\color{ManualGray}内部使用｜技术辅助系统}
\fancyfoot[R]{\small\color{ManualGray}第 \thepage\ 页 / 共 \pageref*{LastPage} 页}
\renewcommand{\headrulewidth}{0.4pt}
\renewcommand{\footrulewidth}{0.4pt}

\fancypagestyle{plain}{
  \fancyhf{}
  \fancyfoot[L]{\small\color{ManualGray}内部使用｜技术辅助系统}
  \fancyfoot[R]{\small\color{ManualGray}第 \thepage\ 页 / 共 \pageref*{LastPage} 页}
  \renewcommand{\headrulewidth}{0pt}
  \renewcommand{\footrulewidth}{0.4pt}
}

\newtcolorbox{importantbox}{
  enhanced,
  breakable,
  colback=ManualLight,
  colframe=ManualBlue,
  boxrule=0.8pt,
  arc=1.5mm,
  left=3mm,right=3mm,top=2mm,bottom=2mm,
  title={重要说明},
  fonttitle=\bfseries
}
\newtcolorbox{warningbox}{
  enhanced,
  breakable,
  colback=orange!6,
  colframe=ManualOrange,
  boxrule=0.8pt,
  arc=1.5mm,
  left=3mm,right=3mm,top=2mm,bottom=2mm,
  title={注意},
  fonttitle=\bfseries
}
\newtcolorbox{dangerbox}{
  enhanced,
  breakable,
  colback=red!4,
  colframe=ManualRed,
  boxrule=0.8pt,
  arc=1.5mm,
  left=3mm,right=3mm,top=2mm,bottom=2mm,
  title={禁止事项},
  fonttitle=\bfseries
}
\newtcolorbox{tipbox}{
  enhanced,
  breakable,
  colback=green!4,
  colframe=ManualTeal,
  boxrule=0.8pt,
  arc=1.5mm,
  left=3mm,right=3mm,top=2mm,bottom=2mm,
  title={操作提示},
  fonttitle=\bfseries
}

\newcommand{\ui}[1]{\textsf{\bfseries #1}}
\newcommand{\code}[1]{\texttt{\small #1}}
\newcommand{\status}[1]{\textsf{#1}}
\newcommand{\blankline}{\par\noindent\rule{\textwidth}{0.4pt}\par}

\newcommand{\MakeManualCover}[5]{
  \begin{titlepage}
    \thispagestyle{empty}
    \vspace*{18mm}
    {\sffamily\bfseries\fontsize{18pt}{22pt}\selectfont\color{ManualTeal} MineGuard 煤炭智能辅助监管系统\par}
    \vspace{8mm}
    \noindent{\color{ManualBlue}\rule{\textwidth}{2.2pt}}
    \vspace{20mm}
    {\centering
      {\sffamily\bfseries\fontsize{31pt}{39pt}\selectfont\color{ManualBlue} #1\par}
      \vspace{9mm}
      {\sffamily\fontsize{16pt}{22pt}\selectfont\color{ManualGray} #2\par}
    }
    \vfill
    \begin{tcolorbox}[colback=ManualLight,colframe=ManualLine,boxrule=0.5pt,arc=1mm]
      \begin{tabularx}{\linewidth}{P{31mm}Y}
        \textbf{适用软件} & #3\\
        \textbf{文档版本} & V1.0\\
        \textbf{发布日期} & 2026 年 8 月 10 日\\
        \textbf{适用对象} & #4\\
        \textbf{文档属性} & 内部使用，正式软件操作与培训材料\\
      \end{tabularx}
    \end{tcolorbox}
    \vspace{8mm}
    {\small\color{ManualGray}#5\par}
    \vspace{12mm}
    {\centering\small\color{ManualGray}MineGuard 项目组\par}
  \end{titlepage}
}

\newcommand{\DocumentControl}[2]{
  \chapter*{文档控制}
  \addcontentsline{toc}{chapter}{文档控制}
  \begin{longtable}{P{34mm}P{110mm}}
    \toprule
    \textbf{项目} & \textbf{内容}\\
    \midrule
    文档名称 & \ManualTitle\\
    软件版本 & #1\\
    文档版本 & V1.0\\
    发布日期 & 2026 年 8 月 10 日\\
    适用范围 & #2\\
    维护原则 & 软件页面、权限或业务流程发生变化时，应同步修订本手册并重新导出 PDF。\\
    \bottomrule
  \end{longtable}

  \section*{修订记录}
  \begin{longtable}{P{20mm}P{30mm}P{25mm}P{70mm}}
    \toprule
    \textbf{版本} & \textbf{日期} & \textbf{状态} & \textbf{说明}\\
    \midrule
    V1.0 & 2026-08-10 & 十量 V3 正式版 & 按当前十量 V3 实现整理角色权限、标准操作、异常处理、安全边界与部署访问说明。\\
    \bottomrule
  \end{longtable}
}


\begin{document}
\MakeManualCover
  {企业端软件使用手册}
  {企业可信数据填报智能体}
  {Enterprise Reporting Agent 0.1.0}
  {企业领导、填报经办人、复核确认人、报送人员和企业系统管理员}
  {本手册用于指导企业侧可信数据整理、复核、确认和报送。软件输出仅作煤炭业务辅助，不替代企业内控、原始凭证审核或监管认定。}

\DocumentControl
  {Enterprise Reporting Agent 0.1.0}
  {当前仓库 \code{agent/} 企业端网页、命令行和独立 HTTP 服务。}

\begin{importantbox}
企业端负责本企业数据的整理、来源说明、人工复核和报送，不查看其他企业数据，也不
作辖区排名或监管定案。企业端与监管平台不共享代码和数据库，双方只通过版本化
HTTP/JSON/HMAC 接口交换报送包与回执。
\end{importantbox}

\tableofcontents
\clearpage

\chapter{系统概述}

\section{系统用途}

企业可信数据填报智能体用于把分散在生产、运销、洗选、库存、计量和安全资料中的
数据整理成结构化草稿，并在正式报送前完成来源追溯、缺项提示、确定性预检、人工
逐项核对和版本确认。系统可以调用煤炭专用工具或已配置的大模型提供候选建议，但
不会替代用户完成企业确认、签名或监管提交。

\section{适用业务}

\begin{itemize}
  \item 建立某矿井、某统计窗口的企业报送草稿；
  \item 导入 JSON、CSV、文字材料或监管事件快照；
  \item 记录生产量、运输量、入洗量、销售量、库存变化及其来源；
  \item 运行完整性、单位、时间、来源、物理关系和历史趋势检查；
  \item 通过煤炭智能任务和煤炭业务对话辅助分析；
  \item 由实名人员逐项复核、确认当前版本并提交监管平台；
  \item 保存平台回执、修订历史和审计记录。
\end{itemize}

\section{两端边界与数据流}

\begin{center}
\begin{tabularx}{0.94\textwidth}{Y c Y}
\textbf{企业原始系统/材料} & $\longrightarrow$ & \textbf{企业填报智能体}\\
生产、计量、洗选、运销、库存、事件材料
& & 草稿、来源摘要、核对、确认、提交记录\\[2mm]
& \textbf{版本化 API + HMAC} & $\downarrow$\\[2mm]
& & \textbf{监管平台}\\
& & 独立验签、交叉验证、时序分析与人工核查
\end{tabularx}
\end{center}

\begin{warningbox}
「平台已接收」仅表示本次报文通过技术接入检查，不表示数据真实、正常、合法、合规，
也不表示监管部门已经审核通过。
\end{warningbox}

\section{企业数据保存在哪里}

企业端默认使用 SQLite 文件保存草稿、来源摘要、人工核对、智能任务、对话、提交
记录和审计信息。开发环境默认位置为：
\begin{quote}
\path{agent/data/enterprise-agent.db}
\end{quote}
生产环境应通过 \path{ENTERPRISE_AGENT_DB} 指定固定的绝对路径，例如：
\begin{quote}
\path{/var/lib/enterprise-agent/enterprise-agent.db}
\end{quote}

推荐每个独立企业部署一套企业端服务和一份独立数据库。同一企业的多个矿井可以在
一套服务内用不同 \code{mine\_id} 管理。企业数据库不应放在监管服务器中，也不应
由多个相互独立企业共用。

\chapter{账号、登录与权限}

\section{登录}

在浏览器打开企业端地址，输入管理员分配的账号和密码。登录后先检查页面显示的
姓名、岗位和权限提示。岗位名称用于展示和留痕，真正可执行的操作由
\code{read}、\code{write}、\code{confirm}、\code{submit} 四项权限决定。

\begin{longtable}{P{25mm}P{42mm}P{78mm}}
\toprule
\textbf{权限} & \textbf{业务含义} & \textbf{主要能力}\\
\midrule
\code{read} & 查看与分析 & 查看草稿、来源、问题、预检、回执和审计；运行只读工具与煤炭对话。\\
\code{write} & 填报与修改 & 新建、编辑、导入、移除未提交草稿；采纳候选建议和批准草稿补丁。\\
\code{confirm} & 复核确认 & 以当前本人身份逐项核对观测，并确认当前修订版本。\\
\code{submit} & 正式报送 & 提交已经有效确认且通过门禁的当前版本，并保存平台回执。\\
\bottomrule
\end{longtable}

\section{推荐岗位分工}

\begin{longtable}{P{31mm}P{38mm}P{77mm}}
\toprule
\textbf{岗位} & \textbf{推荐权限} & \textbf{职责}\\
\midrule
企业查看/领导 & \code{read} & 查看填报进度、异常、回执和审计；不修改数据。\\
填报经办人 & \code{read, write} & 建草稿、导资料、回答缺项、修正数据并预检。\\
复核确认人 & \code{read, confirm} & 对照原始资料逐项核对并确认当前版本。\\
报送人员 & \code{read, submit} & 检查确认和门禁，执行提交并留存回执。\\
全流程人员 & 四项权限 & 技术上可完成全流程；组织制度仍可要求岗位分离。\\
\bottomrule
\end{longtable}

\section{本机演示账号}

未配置正式用户且服务只监听 \code{127.0.0.1} 时，系统可启用本机演示账号：

\begin{center}
\begin{tabular}{ll}
\toprule
账号 & \code{demo}\\
默认密码 & \code{123123123}\\
有效权限 & \code{read, write}\\
限制 & 必须修改密码；禁止正式逐项核对、确认和提交\\
\bottomrule
\end{tabular}
\end{center}

\begin{dangerbox}
演示账号和默认密码不得用于内网生产。企业端当前没有网页自助改密功能；正式账号
由管理员生成密码摘要、更新用户配置并重启服务。不得把明文密码写入仓库、用户 JSON
或启动脚本。
\end{dangerbox}

\section{权限不足时如何处理}

按钮禁用并不一定是故障。先查看页面给出的具体原因：缺少权限、账号待换密、草稿
已经提交、当前版本尚未逐项核对、存在预检阻断、平台未配置或同版本正在提交。
不要通过更换岗位名称、直接改库或要求 AI 绕过门禁。

\chapter{页面与标准流程}

\section{快捷填报与专业工具}

企业端登录后默认进入\ui{快捷填报}。该模式面向日常经办和非技术用户，只突出
当前最需要完成的一项任务、完成进度和下一步按钮；完整的数据校验、权限门禁、
修订留痕和六步流程仍在后台生效，不会因界面简化而跳过。

\begin{itemize}
  \item 有未完成草稿时，首页主按钮显示\ui{继续未完成填报}，优先续办原任务；
  \item 确实需要并行填报另一统计期时，再使用\ui{新建另一份}，避免重复建单；
  \item \ui{当前任务}卡片会按草稿状态提示导入、核对、预检、确认或提交；
  \item 没有相应权限时，按钮只提供查看和定位，不会代替经办人、确认人或提交人执行操作；
  \item 页面顶部的系统状态会说明当前能否编辑、确认或提交及其原因。
\end{itemize}

需要批量导入、智能任务、煤炭业务对话、审计或诊断工具时，可切换到
\ui{专业工具}。切换模式不会创建新草稿、不会保存业务数据，也不会改变当前账号
权限；正在编辑的草稿和页面位置会保留。完成专业操作后可随时返回
\ui{快捷填报}。

\begin{tipbox}
日常填报建议始终跟随\ui{当前任务}按钮操作。只有在核查来源、处理复杂异常或由
技术人员排障时才进入\ui{专业工具}，这样可以减少误操作和无关信息干扰。
\end{tipbox}

\section{六步工作区}

企业端围绕一个草稿组织操作，页面主流程为：

\begin{enumerate}
  \item \ui{基本信息}：企业、矿井、统计窗口、配置版本和运行工况；
  \item \ui{导入来源}：JSON、CSV、文字材料或事件快照；
  \item \ui{提取核对}：查看候选字段、来源位置和可信度，人工采纳或修改；
  \item \ui{提交前预检}：检查缺项、逻辑、来源、单位、时间和业务关系；
  \item \ui{账号确认}：由确认人本人逐项核对并确认当前修订；
  \item \ui{提交回执}：提交监管平台，查看技术接收状态和回执。
\end{enumerate}

\begin{tipbox}
页面保存和确认都绑定「修订号」。出现「草稿已更新」时，应刷新页面、核对差异并重新
执行受影响步骤，不得直接覆盖其他人的新版本。
\end{tipbox}

\section{开始填报前的准备}

\begin{itemize}
  \item 确认企业编号、统一社会信用代码、矿井编号和名称；
  \item 明确统计窗口的开始、结束时间和时区；
  \item 准备经批准的配置 \code{profile\_id + version}；
  \item 准备原始台账、设备导出文件、审批材料和来源编号；
  \item 明确本期工况：生产制度、班次、季节、检修状态和特殊事件；
  \item 确认填报、复核、报送人员使用各自实名账号。
\end{itemize}

\section{步骤一：新建草稿}

\begin{enumerate}
  \item 以具有 \code{write} 权限的账号登录；有未完成任务时点击
  \ui{继续未完成填报}，没有时点击\ui{新建填报}。只有确需并行填报时才点击
  \ui{新建另一份}。
  \item 录入企业和矿井信息，核对编号是否与监管端登记一致。
  \item 设置统计窗口。结束时间必须晚于开始时间，且不得用模糊文字代替。
  \item 选择或填写配置编号及版本；不要把测试配置用于正式报送。
  \item 填写四轴工况及事件说明，并导入与本矿井、统计窗口一致的监管事件快照；
  即使查询结果为空，也必须保留有效的空结果快照。停产、检修、调拨等特殊工况
  应有材料依据。
  \item 保存后记录草稿编号和当前修订号。
\end{enumerate}

\section{步骤二：导入来源材料}

\subsection{JSON 和 CSV}

结构化文件优先走确定性导入。导入后检查列名、单位、时间、指标代码、来源编号、
小数点和空值语义。缺失值必须保持缺失，不能为通过检查而补零。

\subsection{文字材料与模型提取}

已配置模型时，系统可以从用户主动选择或粘贴的文字中生成候选字段。候选值不会自动
写入草稿；经办人必须查看原文位置、字段、单位和可信度后逐项采纳。没有配置模型时，
JSON/CSV 导入、人工编辑、预检、确认和提交仍可正常使用。

\subsection{监管事件快照}

监管事件代码应使用专用快照导入。手工填写空数组不能冒充权威查询结果。报送时平台
会核对快照编号、矿井、时间窗、完整代码集和来源证据摘要。

\section{步骤三：核对候选数据}

对每个观测至少核对：

\begin{itemize}
  \item 指标名称和指标代码是否对应；
  \item 数值、单位、统计口径和正负方向是否正确；
  \item 时间窗口是否完整覆盖本次报送期；
  \item 来源编号、来源系统、文件名及原始位置是否可追溯；
  \item 是否存在重复导入、累计表复位、跨期库存或缺测；
  \item 来源真实性声明是否与实际材料一致。
\end{itemize}

人工修改受签字段会改变内容摘要，使原来源签名失效。此时应从来源系统重新导出并由
来源网关重新签发，企业填报智能体不得持有来源密钥或替人工数据补签。

\section{步骤四：运行预检}

点击 \ui{提交前预检}，按严重程度处理：

\begin{longtable}{P{30mm}P{44mm}P{72mm}}
\toprule
\textbf{类型} & \textbf{含义} & \textbf{处理}\\
\midrule
阻断 & 不能进入确认或提交 & 修正数据、来源、权限或配置后重新预检。\\
需说明 & 存在异常或口径不清 & 结合原始资料补充说明，必要时退回来源部门核验。\\
提示 & 不阻断但值得关注 & 确认已知并保留说明，继续人工复核。\\
通过 & 自动规则未发现阻断 & 仍须人工核对；不代表事实真实或监管认可。\\
\bottomrule
\end{longtable}

\section{步骤五：逐项复核和人工确认}

\begin{enumerate}
  \item 确认人用本人账号登录，核对草稿编号、矿井、窗口和当前修订号。
  \item 对照原始台账、设备记录和审批材料，逐条点击核对并保存本人核对记录。
  \item 确认进度为「当前本人全部已核对」。其他账号的记录不能替代本人复核。
  \item 再次运行预检，确保阻断项为零。
  \item 阅读真实性声明，主动勾选接受，点击 \ui{完成人工确认}。
  \item 检查页面显示确认人、岗位、确认时间、内容摘要和有效修订号。
\end{enumerate}

\begin{warningbox}
确认绑定当前修订号和内容摘要。确认后任何业务内容变化都会使旧确认失效；受影响
观测需要重新逐项核对，随后重新预检和确认。
\end{warningbox}

\section{步骤六：提交与保存回执}

\begin{enumerate}
  \item 报送人员检查当前版本已有有效确认，账号具备 \code{submit} 权限。
  \item 检查平台地址、客户端编号、运输密钥和契约版本已经配置。
  \item 检查来源签名、事件快照和确认人登记均满足平台门禁。
  \item 点击提交并等待明确结果。处理中不要刷新后反复生成新提交。
  \item 如页面超时，先刷新提交记录，核对服务端是否已经收到。
  \item 成功后记录回执号、接收时间、报送摘要和草稿修订号。
\end{enumerate}

网络故障重试应沿用原幂等键；内容、权限、配置或治理错误必须先修正原因。成功报送
的草稿保持锁定，更正应建立新报送并保留原记录。

\chapter{煤炭智能助手}

\section{一键体检和智能任务}

系统提供煤流、库存、单位、时间、来源血缘、物理复算、历史趋势、煤质场景等只读
确定性工具。相同输入可复算并不等于输入真实。智能任务可以规划工具调用和形成解释，
但不能执行企业确认、来源签名或监管提交。

\begin{longtable}{P{38mm}P{108mm}}
\toprule
\textbf{功能} & \textbf{正确用法}\\
\midrule
一键煤炭体检 & 绑定当前草稿，快速查看完整性、物理关系、历史和来源问题。\\
确定性工具模式 & 不调用第三方模型，适合保密场景和可复算核验。\\
自动规划模式 & 由已配置模型选择只读工具；关键结论仍以工具证据为准。\\
草稿补丁 & 只有 \code{write} 用户可逐项查看字段、前后值和参数摘要后批准。\\
任务取消/拒绝 & 创建人可停止本人任务；已发生的审计记录继续保留。\\
\bottomrule
\end{longtable}

\section{煤炭业务对话}

\begin{itemize}
  \item 煤炭通识：可解释煤质、燃烧、生产、洗选、计量、库存和安全常识；
  \item 近期煤炭新闻：使用已配置的联网搜索获取来源片段，再由 AI 归纳；
  \item 企业实际数据：必须绑定正确草稿，以只读工具证据回答；
  \item 非煤炭、越界操作、提示注入或绕过审批：系统拒绝。
\end{itemize}

新闻检索失败时，系统不会用模型记忆伪装成近期新闻。用户应打开来源核对日期、主体
和全文。对话答案、搜索摘要和通识说明均不是报送凭证。

\section{模型 API 配置}

模型配置通过环境变量在服务启动时读取，修改后需要重启。示例变量以部署模板和
\code{agent/.env.example} 为准，常用项包括：

\begin{verbatim}
DEEPSEEK_API_KEY=由秘密管理工具提供
DEEPSEEK_BASE_URL=https://api.deepseek.com
DEEPSEEK_MODEL=管理员批准的模型名称
\end{verbatim}

\code{.env} 文件不会自动加载；可由 systemd 的 \code{EnvironmentFile}、容器秘密
或受控启动环境注入。API Key 不得写入网页、对话、草稿、代码仓库或日志。

\begin{dangerbox}
不得向 AI 粘贴密码、Cookie、CSRF 令牌、API Key、Bearer Token、运输 HMAC、
来源签名密钥、私钥、完整监管报文或无关个人敏感信息。
\end{dangerbox}

\chapter{草稿生命周期与审计}

\section{草稿状态}

\begin{longtable}{P{34mm}P{112mm}}
\toprule
\textbf{状态} & \textbf{说明}\\
\midrule
编辑中 & 可由 \code{write} 用户修改；每次业务变更产生新修订。\\
待复核 & 经办完成，等待确认人逐项核对。\\
已确认 & 当前修订已由本人逐项核对并确认；修改会使确认失效。\\
提交中 & 正在请求监管平台；不得删除或并发修改。\\
提交成功 & 已取得技术接收回执；原草稿锁定。\\
提交失败/未知 & 先根据原幂等键和服务端记录核对，再决定重试或修正。\\
已移除 & 未提交草稿的软删除状态；审计仍保留。\\
\bottomrule
\end{longtable}

\section{删除与更正}

具有 \code{write} 权限的用户可以在二次文字确认后移除「未提交且没有提交处理中记录」
的草稿。该操作为软删除，普通列表不再展示，但审计记录不抹除。正在提交或已经成功
提交的草稿不能删除。正式更正应新建更正报送，引用原回执并说明原因。

\section{审计与完整性}

系统保留修订、来源、候选采纳、核对、确认、任务、工具批准、提交和回执事件。
哈希链可以发现意外损坏或局部改写，但不是外部数字签名、可信时间戳或不可篡改存储。
如发现完整性失败，应停止批准、确认和提交，保全数据库副本并进入安全事件处置。

\chapter{常见问题与故障处理}

\begin{longtable}{P{45mm}P{101mm}}
\toprule
\textbf{现象} & \textbf{处理办法}\\
\midrule
终端显示启动信息后不动 & 服务正在等待浏览器请求，不是卡死。保持终端运行，浏览器访问配置地址；停止时按 Ctrl+C。\\
\code{enterprise-agent: command not found} & 进入 \code{agent/}，激活安装该包的虚拟环境并执行 \code{pip install -e .}；也可临时使用 \code{PYTHONPATH=src python -m enterprise\_agent ...}。\\
\code{PYTHONPATH=src: command not found} & 不要粘贴终端的控制字符；手工输入完整命令，等号两侧不要加空格。\\
登录成功但按钮禁用 & 查看实际权限、待换密标记、草稿状态和门禁原因。演示账号不能确认或提交。\\
模型 API 已配置但新闻不可用 & 模型 API 和联网搜索是两项独立能力；检查搜索配置、网络和来源返回。\\
修改后确认消失 & 正常的版本保护。重新核对受影响观测、预检并确认新修订。\\
提交超时 & 先刷新提交记录和平台回执，禁止盲目创建新幂等键重复提交。\\
重启后看不到原数据 & 核对 \code{ENTERPRISE\_AGENT\_DB} 是否指向原绝对路径，避免进程打开另一份相对路径数据库。\\
来源签名失效 & 从原来源系统重新导出并由来源网关签发，不能在填报端补签。\\
\bottomrule
\end{longtable}

\chapter{启动、访问与生产部署}

\section{本机快速启动}

\begin{verbatim}
cd /home/sevan/coral/agent
python3 -m venv .venv
. .venv/bin/activate
python -m pip install -e .
enterprise-agent serve --host 127.0.0.1 --port 8090
\end{verbatim}

浏览器打开：
\begin{quote}
\url{http://127.0.0.1:8090/}
\end{quote}
健康检查为：
\begin{quote}
\url{http://127.0.0.1:8090/api/v1/health}
\end{quote}

\section{正式服务}

生产建议使用独立低权限账号、固定代码目录、固定状态目录、systemd 和 HTTPS
反向代理。仓库提供：

\begin{itemize}
  \item \code{agent/deploy/enterprise-agent.env.example}
  \item \code{agent/deploy/enterprise-agent.service.example}
  \item \code{agent/deploy/nginx-enterprise-agent.conf.example}
\end{itemize}

后端只监听 \code{127.0.0.1:8090}，由 Nginx 对内提供 HTTPS 443。不要直接把
8090 暴露到局域网或公网。正式环境设置唯一
\path{ENTERPRISE_AGENT_PUBLIC_ORIGIN}，通过 HTTPS 访问时启用 Secure Cookie。

\section{监管平台连接参数}

企业端至少需要与监管端对齐以下配置：

\begin{itemize}
  \item \code{PLATFORM\_BASE\_URL}：监管平台 HTTPS 地址；
  \item \code{PLATFORM\_CLIENT\_ID}：监管端登记的企业客户端编号；
  \item \code{PLATFORM\_TRANSPORT\_HMAC\_SECRET}：运输 HMAC 密钥；
  \item 企业编号和允许报送的矿井编号范围；
  \item 契约版本、配置版本、来源注册、事件快照和确认人登记。
\end{itemize}

运输密钥只保护企业端到监管平台的报文，不等同于设备来源签名密钥，二者严禁复用。

\section{备份}

SQLite 运行时可能同时存在 \code{-wal} 和 \code{-shm} 文件。不得在服务运行中只
复制单个 \code{.db} 文件。应在维护窗口停止服务后备份整个状态目录，或使用经验证
的 SQLite 在线备份工具。恢复后核对草稿、回执、提交锁定、审计链和任务事件链。
状态盘及备份介质应使用单位认可的静态加密和访问控制。

\chapter{岗位操作清单}

\section{经办人交接清单}

\begin{itemize}
  \item[$\square$] 企业、矿井、统计窗口和配置版本正确；
  \item[$\square$] 工况和特殊事件有材料依据；
  \item[$\square$] 每项观测的来源、单位、时间和口径可追溯；
  \item[$\square$] 缺失值未补零，累计表复位和库存跨期已说明；
  \item[$\square$] 预检阻断项为零，需说明项已处理；
  \item[$\square$] 已向确认人交接草稿编号和当前修订号。
\end{itemize}

\section{确认人清单}

\begin{itemize}
  \item[$\square$] 使用本人实名账号；
  \item[$\square$] 核对正确矿井、窗口和修订号；
  \item[$\square$] 本人逐项核对全部观测；
  \item[$\square$] 预检阻断项为零；
  \item[$\square$] 阅读并主动接受真实性声明；
  \item[$\square$] 页面显示当前修订已有效确认。
\end{itemize}

\section{报送人员清单}

\begin{itemize}
  \item[$\square$] 当前版本的人工确认有效；
  \item[$\square$] 平台配置、来源签名和治理门禁有效；
  \item[$\square$] 同版本没有正在进行的提交；
  \item[$\square$] 提交后保存回执号、接收时间和摘要；
  \item[$\square$] 明确「技术接收」不等于「监管通过」。
\end{itemize}

\appendix
\chapter{术语速查}

\begin{longtable}{P{38mm}P{108mm}}
\toprule
\textbf{术语} & \textbf{含义}\\
\midrule
草稿修订号 & 每次业务内容变化后的版本编号，用于防止并发覆盖。\\
内容摘要 & 对规范化报送内容计算的摘要，用于绑定确认、签名和幂等。\\
来源凭证 & 说明数据来自何处、何时、何位置以及是否有来源签名。\\
来源签名 & 由设备或采集网关使用来源密钥生成；企业填报端不得补签。\\
运输 HMAC & 企业端与监管平台之间的接口认证和报文完整性保护。\\
人工确认 & 当前实名账号对当前修订执行的主动点击确认；不自动等同于合格电子签章。\\
幂等键 & 网络重试时识别「同一提交」的标识，防止重复入库。\\
技术接收 & 平台接入成功；尚未形成监管分析或业务结论。\\
\bottomrule
\end{longtable}

\chapter{责任边界}

\begin{importantbox}
本系统是煤炭数据填报与技术核查辅助工具。任何模型回答、工具结果、预检通过、风险
提示或平台回执，都不能替代原始凭证审核、现场核验、专业判断、企业内控程序和有权
机关依法作出的认定。发生安全生产紧急情况时，应执行本单位现行应急预案和适用规程，
不得等待本系统或 AI 给出指令。
\end{importantbox}

\end{document}
